\section{Funções Escolha}\label{sec:funcoesEscolha}
\index{função escolha|see{função de escolha}}\index{função de escolha}

Frequentemente, em matemática, temos uma família indexada $(A_i: i\in I)$ de conjuntos não vazios e queremos ``fixar um elemento $a_i \in A_i$ para cada $i\in I$''.
A Teoria dos Conjuntos, com os axiomas que temos até então, não é suficiente para garantir a existência de uma tal família $(a_i: i\in I)$ em qualquer tal situação.

Nesta seção, vamos estudar tais famílias.

\begin{definition}[$Z^--P$]\label{def:funcaoEscolha}\index{função de escolha}
    Seja $\mathcal C$ uma coleção de conjuntos não vazios.
    Uma \emph{função de escolha} para $\mathcal C$ é uma função $c$ de domínio $\mathcal C$ tal que, para cada $a\in \mathcal C$, $c(a)\in a$.
\end{definition}
\begin{example}[$Z^--P$]
    Seja $\mathcal C=\{\{1, 2\}, \{3, 4\}\}$.

    Temos que $c:\mathcal C\to \{1, 2, 3, 4\}$ dada por

    \[
        c(a) = \begin{cases}
            1, & \text{se } a=\{1, 2\}\\
            3, & \text{se } a=\{3, 4\}\\
        \end{cases}
    \]
    é uma função de escolha para $\mathcal C$.
\end{example}
Na Seção~\ref{sec:relacoesDeEquivalencia}, foi definida a noção de partição.
Lembremos que, nesse texto, uma partição de um conjunto $X$ é uma coleção $P$ de subconjuntos não vazios de $X$ dois-a-dois disjuntos e cuja união é $X$.

\begin{definition}[$Z^--P$]
    Seja $X$ um conjunto e $P$ uma partição de $X$.
    Um \emph{conjunto transversal} para $P$, ou \emph{conjunto de representantes} para $P$, é um subconjunto $T$ de $X$ tal que, para cada $a\in P$, $T\cap a$ é um conjunto unitário.
\end{definition}

\begin{example}[$Z^--P$]
    Seja $X=\{1, 2, 3, 4\}$ e $P=\{\{1, 2\}, \{3, 4\}\}$.
    Temos que $T=\{1, 3\}$ é um conjunto transversal para $P$, pois $T\cap \{1, 2\}=\{1\}$ e $T\cap \{3, 4\}=\{3\}$ são conjuntos unitários.
\end{example}


\begin{definition}[$Z^--P$]
    Seja $(A_i: i\in I)$ uma família indexada de conjuntos não vazios.
    Uma \emph{função de escolha indexada} para $(A_i: i\in I)$ é uma função $c$ de domínio $I$ tal que, para cada $i\in I$, $c(i)\in A_i$.
\end{definition}
Novamente, notemos que se $A_i$ é vazio, não há nenhuma restrição sobre o valor de $c(i)$.

\begin{example}[$Z^--P$]
    Seja $I=\{1, 2\}$, $A_1=\{1, 2\}$ e $A_2=\{3, 4\}$.

    Temos que $c:I\to \{1, 2, 3, 4\}$ dada por
    \[
        c(i) = \begin{cases}
            2, & \text{se } i=1\\
            3, & \text{se } i=2\\
        \end{cases}
    \]
    é uma função de escolha indexada para $(A_i: i\in I)$.
\end{example}

Nesse contexto, temos:

\begin{proposition}[$ZF^--P$]\label{prop:equivalenciaAxiomaEscolha}
    São equivalentes.
    \begin{enumerate}[label=(\alph*)]
        \item Para todo conjunto $X$ e toda partição $P$ de $X$, existe um conjunto transversal para $P$.
        \item Para toda família indexada $(A_i: i\in I)$ de conjuntos não vazios, existe uma função de escolha indexada para $(A_i: i\in I)$.
        \item Para toda coleção de conjuntos não vazios $\mathcal C$, existe uma função de escolha para $\mathcal C$.
    \end{enumerate}
\end{proposition}
\begin{proof}
    (a) $\Rightarrow$ (b): seja $(A_i: i\in I)$ uma família indexada de conjuntos não vazios.
    Considere $X=\bigcup_{i\in I} \{i\}\times A_i$.
    Seja $T\subseteq X$ tal que para cada $i\in I$, $T\cap (\{i\}\times A_i)$ é um conjunto unitário.
    Então $T:I\rightarrow \bigcup_{i\in I} A_i$ é uma função de escolha.

    (b) $\Rightarrow$ (c): seja $\mathcal C$ uma coleção de conjuntos não vazios.
    Seja $I=\mathcal C$ e considere a família $(A_a: a\in I)$ dada por $A_a=a$.
    Seja $f:\mathcal C\to \bigcup_{a\in \mathcal C} a$ uma função de escolha indexada para $(A_a: a\in I)$.
    Segue que, para cada $a\in \mathcal C$, $f(a)\in A_a=a$, de modo que $f$ é uma função de escolha para $\mathcal C$.

    (c) $\Rightarrow$ (a): seja $X$ um conjunto e $P$ uma partição de $X$.
    Seja $f:P\to X$ uma função de escolha para $P$.
    Então $T=f[P]$ é como desejado.
\end{proof}

Observação: com o axioma da potência, o produto cartesiano $\prod_{i\in I} A_i$ está bem definido.
Com isso, na presença do Axioma da Potência, o item (b) da proposição anterior é equivalente a dizer que, para toda família indexada $(A_i: i\in I)$ de conjuntos não vazios, $\prod_{i\in I} A_i \neq \emptyset$.

Neste texto, tomaremos, pela sua maior simplicidade estrutural, o primeiro item da proposição anterior como a definição do Axioma da Escolha.
Para tornar seu enunciado ainda mais simples, omitiremos o conjunto $X$, que pode ser recuperado, se desejado, unindo-se os conjuntos da partição.
\begin{axiom}[Axioma da Escolha]\label{ax:escolhaParticao}\index{axioma!da escolha}
    Para toda partição $P$, existe um conjunto $T$ que intersecta cada elemento de $P$ em um conjunto unitário.
    Em símbolos:

    \[
    \forall P\Bigl(\forall a, b \in P(a\neq b \Rightarrow \nexists x(x\in a \wedge x\in b)) \rightarrow \exists T\,\forall a\in P\,\exists! x(x\in T \wedge x\in a)\Bigr).
    \]
\end{axiom}

É imediato que, na presença de $ZF^--P$, o Axioma da Escolha é equivalente às cláusulas da Proposição~\ref{prop:equivalenciaAxiomaEscolha}.
O mesmo é válido em $Z^-$.
